\documentclass[12pt]{article}
\usepackage[utf8]{inputenc}
\usepackage[T1]{fontenc}
\usepackage{hyperref}
\title{Shekhawati Language Resources for K-12 Curriculum Development}
\author{Compiled by Claude Sonnet 5.5 (claude-sonnet-5-5) for Andr\'as Kornai}
\date{}
\begin{document}
\maketitle

\section{Introduction}
Shekhawati (ISO 639-3 swv; Glottocode shek1243) is an Indo-European (Indo-Aryan) language of the Shekhawati region of Rajasthan, India. Wikipedia estimates about three million speakers. Some sources treat it as a dialect of Rajasthani or Marwari, so searches may need those broader terms. The searches found one dedicated descriptive grammar, no dedicated dictionary, and no school materials.

\section{Primary Educational Resources}
No primers, textbooks or school programmes specific to Shekhawati were found in search results.

\section{Linguistic and Reference Materials}
\begin{itemize}
\item Gusain, Lakhan (2001), \emph{Shekhawati}, Languages of the World/Materials 385, Lincom Europa (about 90 pages; phonology, morphology, syntax, short interlinear text). Review notice: \url{https://linguistlist.org/issues/12/2004/}
\item A second LINGUIST List notice mentioning the same series: \url{https://linguistlist.org/issues/17/3764/}
\item Columbia University AMESA scholar profile for Lakhan Gusain: \url{https://amesa.library.columbia.edu/scholars/718?locale=en}
\item Wikipedia, ``Shekhawati language'' (thin and partly uncited): \url{https://en.wikipedia.org/wiki/Shekhawati_language}
\item ASJP Shekhawati wordlist: \url{https://asjp.clld.org/languages/SHEKHAWATI}
\item Grierson's \emph{Linguistic Survey of India} (1908), Vol.\ IX Part II, which has a section on Shekhawati; the search result is the BnF catalogue record: \url{https://catalogue.bnf.fr/ark:/12148/cb160322897}
\end{itemize}

\section{Religious and Community Resources}
None found in search results.

\section{Digital and Audio Resources}
Only the ASJP wordlist above, which is a small lexical dataset. No audio or corpora were found.

\section{Curriculum Development Context}
The language is not known to be taught in schools. Rajasthani and Hindi are the likely school-language context, but the searches did not confirm this. Gusain's grammar and Grierson's survey could inform orthography and grammar. Classroom materials would have to be created.

\end{document}